\chapter{Bose 统计和 Fermi 统计}

\section{为什么需要量子统计}

当热德布罗意波长与粒子平均间距可比时，
\[
  n\lambda_T^3 \gtrsim 1,
\]
粒子波包明显重叠。此时“给粒子贴标签”的经典图像失效，必须使用全同粒子统计。

全同粒子的多体波函数在交换两个粒子时只有两类可能：
\[
  \Psi(\ldots,i,\ldots,j,\ldots)
  =
  \begin{cases}
    +\Psi(\ldots,j,\ldots,i,\ldots), & \text{Bose},\\
    -\Psi(\ldots,j,\ldots,i,\ldots), & \text{Fermi}.
  \end{cases}
\]
这两个符号最终变成平均占据数中的 $-1$ 和 $+1$。

\section{Bose 与 Fermi 分布的统一形式}

令
\[
  z=e^{\beta\mu}
\]
为逸度。Bose 和 Fermi 的平均占据数为
\[
  \bar n(\epsilon)
  =
  \frac{1}{z^{-1}e^{\beta\epsilon}+\eta}
  =
  \frac{1}{e^{\beta(\epsilon-\mu)}+\eta},
\]
\[
  \eta=
  \begin{cases}
    -1, & \text{Bose},\\
    +1, & \text{Fermi}.
  \end{cases}
\]
Bose 分布允许低能态出现很大的占据数；Fermi 分布始终满足
\[
  0\le \bar n(\epsilon)\le 1.
\]

\section{Bose 气体：多粒子可以进入同一态}

Bose 气体的粒子数为
\[
  N=\sum_l
  \frac{\omega_l}{e^{\beta(\epsilon_l-\mu)}-1}.
\]
为了使分母为正，化学势必须满足
\[
  \mu<\epsilon_0.
\]
若取基态能量 $\epsilon_0=0$，则 $\mu<0$。

温度降低时，激发态最多能容纳的粒子数有限。剩余粒子会集中进入基态，出现 Bose-Einstein 凝聚。三维自由 Bose 气体的凝聚条件可写成数量级关系
\[
  n\lambda_T^3\simeq 2.612.
\]

\begin{keybox}
\textbf{Bose 图像：}不是因为粒子互相吸引才凝聚，而是因为全同玻色子的计数规则允许许多粒子占据同一量子态。
\end{keybox}

\section{光子气体：粒子数不守恒}

光子可以被物质吸收和发射，粒子数不守恒，因此
\[
  \mu=0.
\]
光子的平均占据数为 Planck 分布
\[
  \bar n(\omega)
  =
  \frac{1}{e^{\beta\hbar\omega}-1}.
\]
黑体辐射能量密度满足
\[
  u(T)=aT^4,\qquad P=\frac{u}{3}.
\]
因此总能量和热容标度为
\[
  U\propto VT^4,\qquad C_V\propto VT^3.
\]

\section{Fermi 气体：Pauli 不相容}

Fermi 分布为
\[
  \bar n(\epsilon)
  =
  \frac{1}{e^{\beta(\epsilon-\mu)}+1}.
\]
零温极限下，它变成阶跃函数：
\[
  \bar n(\epsilon)=
  \begin{cases}
    1, & \epsilon<\epsilon_F,\\
    0, & \epsilon>\epsilon_F.
  \end{cases}
\]
也就是说，$T=0$ 时费米子从最低能级开始逐层填满，一直填到 Fermi 能 $\epsilon_F$。

三维自由电子气中
\[
  k_F=(3\pi^2 n)^{1/3},
\qquad
  \epsilon_F=\frac{\hbar^2k_F^2}{2m}.
\]
低温下只有 Fermi 面附近宽度约 $\kb T$ 的粒子能被热激发，所以电子热容不是经典的 $\frac{3}{2}N\kb$，而是
\[
  C_V\propto N\kb \frac{T}{T_F}.
\]

\section{三种统计的概念对比}

\begin{center}
\small
\begin{tabular}{p{0.16\textwidth} p{0.34\textwidth} p{0.40\textwidth}}
\toprule
统计 & 计数规则 & 分布特征与典型现象 \\
\midrule
Boltzmann & 粒子近似可分辨，交换忽略 & $\bar n=e^{-\beta(\epsilon-\mu)}$；经典理想气体 \\
Bose & 全同、对称、可多重占据 & 低能态占据可很大；BEC, 黑体辐射 \\
Fermi & 全同、反对称、每态最多一个 & $0\le\bar n\le1$；Fermi 海，电子简并压 \\
\bottomrule
\end{tabular}
\end{center}

\begin{keybox}
\textbf{复习顺序：}先判断粒子是否全同、是否稀薄；
再判断粒子数是否守恒；最后选择 Boltzmann, Bose 或 Fermi 分布。
\end{keybox}
